\documentclass[10pt,a4paper]{amsart}
\usepackage[margin=19mm]{geometry}
\usepackage{amsmath,amssymb,mathtools}
\usepackage[scr=rsfs,cal=euler]{mathalfa}
\usepackage{xfrac}
\usepackage[unicode,hidelinks,bookmarks=false]{hyperref}
\newcommand{\ie}{i.e.\ }
\newcommand{\cf}{cf.\ }
\newcommand{\resp}{respectively\ }
\newcommand{\Subsection}{\subsection}
\providecommand{\nobreakdash}{\nobreak\hbox{-}}
\setlength{\emergencystretch}{2em}
\multlinegap0pt
\usepackage{amssymb}
\usepackage{xcolor}
\usepackage{tcolorbox}
\usepackage{bm}
\newtheorem{lem}{Lemma}
\title{Beyond Expected Information Gain: Stable Bayesian Optimal Experimental Design with Integral Probability Metrics and Plug-and-Play Extensions}
\author{Di Wu, Ling Liang, Haizhao Yang}
\date{Source: arXiv:2604.21849v1 (CC BY 4.0)}
\begin{document}
\maketitle

\noindent\emph{Source and licence.} Beyond Expected Information Gain: Stable Bayesian Optimal Experimental Design with Integral Probability Metrics and Plug-and-Play Extensions. Version arXiv:2604.21849v1 (CC BY 4.0), distributed by arXiv under the Creative Commons Attribution 4.0 International licence, http://creativecommons.org/licenses/by/4.0/, as recorded in the arXiv licence metadata for this exact version and shown on the abstract page https://arxiv.org/abs/2604.21849v1. Full licence notice: \texttt{licenses/arxiv-2604.21849-CC-BY-4.0.txt}.\par\smallskip

\noindent\textit{Editorial scope.} Counted unit: the article's lem lem:anchor (source0.tex lines 877--887). The article's complete original proof of it is delivered with it (source0.tex lines 1386--1390, one \texttt{proof} environment of 5 lines, which the article places away from the statement). No mathematics is written, repaired or re-derived: the statement and the proof are the article's own lines, in the article's own notation, and no retained line is substituted or re-flowed.  No further source-local context is needed: every cross-reference of every retained line resolves inside the two delivered spans.  Nothing was omitted from any retained span, so the package carries no omission record.  No retained line contains a citation, so the wrapper carries no bibliography and no bibliography entry.  No retained line contains a figure environment, an image-inclusion command or drawing code, so no image is delivered and the package declares no asset dependency.  Editorial formatting only, in addition: an AMS \texttt{amsart} preamble that restates the article's own declarations and macros used by the retained lines, namely the environments \texttt{lem} on separate counters, together with the standard packages those lines need.  The retained environments print in this document as Lemma 1 in order of appearance, whereas the article prints the same items with the numbering of its own full document.  The complete native source of the article as distributed in the arXiv e-print for this exact version is bundled unchanged as \texttt{sources/199051/source0.tex}.
\begin{lem}[Prior-Anchored Supremum] \label{lem:anchor}
    Let $\mathcal{F}$ be a symmetric function class defining an IPM $\gamma_{\mathcal{F}}$ and $\tilde{\mu}$ a perturbed probability measure. The prior-anchored subspace is defined as follows:
    $$\mathcal{F}_{\text{prior}} = \left\{ g \in \mathcal{F} \ \Big| \ \int_{\mathcal{X}} g(x) \tilde{\mu}(dx) = 0 \right\}$$
    Then, the following statements hold:
    \begin{itemize}
        \item The IPM between the two densities can be evaluated exactly via
        $$\sup_{f \in \mathcal{F}} \left| \int_{\mathcal{X}} f(x) p(dx) - \int_{\mathcal{X}} f(x) q(dx) \right| = \sup_{g \in \mathcal{F}_{\text{prior}}} \left| \int_{\mathcal{X}} g(x) p(dx) - \int_{\mathcal{X}} g(x) q(dx) \right|.$$
        \item Specifically, if we set $\mathcal{X}$ to be the parameter space $\mathcal{X}$ and $p=\tilde{\mu}, q=\tilde{\mu}^y$, the IPM between the prior and posterior can be evaluated exactly over this restricted subspace
        $$\gamma_{\mathcal{F}}(\tilde{\mu}, \tilde{\mu}^y) = \sup_{g \in \mathcal{F}_{\text{prior}}} \left| \int_{\mathcal{X}} g(x) \tilde{\mu}^y(dx) \right| \triangleq D(y)$$
    \end{itemize}
\end{lem}

\begin{proof}[Proof of Lemma \ref{lem:anchor}]
    Let $H(f) = \left| \int_{\mathcal{X}} f(x) p(dx) - \int_{\mathcal{X}} f(x) q(dx) \right|$. We wish to prove $\sup_{f \in \mathcal{F}} H(f) = \sup_{g \in \mathcal{F}_{\text{prior}}} H(g)$. Because $\mathcal{F}_{\text{prior}} \subset \mathcal{F}$, it trivially holds that $\sup_{f \in \mathcal{F}} H(f) \ge \sup_{g \in \mathcal{F}_{\text{prior}}} H(g)$. To prove the reverse inequality, consider any arbitrary function $f \in \mathcal{F}$. Let $c = \int_{\mathcal{X}} f(x) p(dx)$. By the shift-invariance of $\mathcal{F}$, the function $g(x) = f(x) - c$ belongs to $\mathcal{F}$. By construction, $\int_{\mathcal{X}} g(x) p(dx) = 0$, meaning that $g \in \mathcal{F}_{\text{prior}}$. We evaluate the posterior discrepancy score $H(g)$ for this shifted function:$$H(g) = \left| \int_{\mathcal{X}} (f(x) - c) p(dx) - \int_{\mathcal{X}} (f(x) - c) q(dx) \right|$$ 
    By the linearity of the integral, we can separate the constant $c$. Since $\mu^y$ and $\tilde{\mu}^y$ are both valid probability measures, they integrate to 1 (i.e., $\int d\mu^y = 1$ and $\int d\tilde{\mu}^y = 1$). The constants $c$ perfectly cancel. Hence, it holds that $$H(g) = \left| \int_{\mathcal{X}} f(x) p(dx) - \int_{\mathcal{X}} f(x) q^y(dx) \right| = H(f)$$
    Because every function $f \in \mathcal{F}$ has a corresponding shifted function $g \in \mathcal{F}_{\text{prior}}$ that yields the exact same absolute difference, the supremum over the full class cannot exceed the supremum over the restricted class. Thus, the supremums are strictly equal.
\end{proof}

\end{document}
